\documentclass[11pt]{article}

\usepackage{amsmath}
\usepackage{amssymb}
\usepackage{amsthm}
\usepackage{mathpartir} % for inference rules; if unavailable, rules are written with \dfrac as fallback
\usepackage{tikz}
\usepackage{hyperref}

\newtheorem{example}{Example}

\title{CS-462 -- Programming Languages Foundations \\ Lecture \#2: Inductive definitions and proofs}
\author{EPFL}
\date{}

\begin{document}

\maketitle

\noindent In this lecture, we will use the semantics of our simple language of arithmetic
expressions,
\[
e ::= x \mid n \mid e_1 + e_2 \mid e_1 * e_2 \mid x := e_1 ; e_2,
\]
to express useful program properties, and we will prove these properties by induction.

\section{Program Properties}

There are a number of interesting questions about a language one can ask: Is it deterministic?
Are there non-terminating programs? What sorts of errors can arise during evaluation? Having a
formal semantics allows us to express these properties precisely.

\begin{itemize}
    \item \textbf{Determinism:} Evaluation is deterministic,
    \[
    \forall e \in \mathbf{Exp}.\ \forall \sigma, \sigma', \sigma'' \in \mathbf{Store}.\ \forall
    e', e'' \in \mathbf{Exp}.
    \]
    \[
    \text{if } \langle \sigma, e \rangle \to \langle \sigma', e' \rangle \text{ and } \langle
    \sigma, e \rangle \to \langle \sigma'', e'' \rangle \text{ then } e' = e'' \text{ and }
    \sigma' = \sigma''.
    \]

    \item \textbf{Termination:} Evaluation of every expression terminates,
    \[
    \forall e \in \mathbf{Exp}.\ \forall \sigma \in \mathbf{Store}.\ \exists \sigma' \in
    \mathbf{Store}.\ \exists e' \in \mathbf{Exp}.\ \langle \sigma, e \rangle \to^* \langle
    \sigma', e' \rangle \text{ and } \langle \sigma', e' \rangle \not\to,
    \]
    where $\langle \sigma', e' \rangle \not\to$ is shorthand for $\neg(\exists \sigma'' \in
    \mathbf{Store}.\ \exists e'' \in \mathbf{Exp}.\ \langle \sigma', e' \rangle \to \langle
    \sigma'', e'' \rangle)$.
\end{itemize}

It is tempting to want the following soundness property,

\begin{itemize}
    \item \textbf{Soundness:} Evaluation of every expression yields an integer,
    \[
    \forall e \in \mathbf{Exp}.\ \forall \sigma \in \mathbf{Store}.\ \exists \sigma' \in
    \mathbf{store}.\ \exists n' \in \mathbf{Int}.\ \langle \sigma, e \rangle \to^* \langle
    \sigma', n' \rangle,
    \]
\end{itemize}

\noindent but unfortunately it does not hold in our language! For example, consider the
totally-undefined function $\sigma$ and the expression $i + j$. The configuration $\langle
\sigma, i + j \rangle$ is \emph{stuck}---it has no possible transitions---but $i + j$ is not an
integer. The problem is that $i + j$ has \emph{free variables} but $\sigma$ does not contain
mappings for those variables.

To fix this problem, we can restrict our attention to \emph{well-formed} configurations
$\langle \sigma, e \rangle$, where $\sigma$ is defined on (at least) the free variables in $e$.
This makes sense as evaluation typically starts with a \emph{closed} expression. We can define
the set of free variables of an expression as follows:
\begin{align*}
fvs(x) &\triangleq \{x\} \\
fvs(n) &\triangleq \{\} \\
fvs(e_1 + e_2) &\triangleq fvs(e_1) \cup fvs(e_2) \\
fvs(e_1 * e_2) &\triangleq fvs(e_1) \cup fvs(e_2) \\
fvs(x := e_1 ; e_2) &\triangleq fvs(e_1) \cup (fvs(e_2) \setminus \{x\})
\end{align*}

Now we can formulate two properties that imply a variant of the soundness property above:

\begin{itemize}
    \item \textbf{Progress:} For each expression $e$ and store $\sigma$ such that the free
    variables of $e$ are contained in the domain of $\sigma$, either $e$ is an integer or there
    exists a possible transition for $\langle \sigma, e \rangle$,
    \[
    \forall e \in \mathbf{Exp}.\ \forall \sigma \in \mathbf{Store}.
    \]
    \[
    fvs(e) \subseteq \operatorname{dom}(\sigma) \implies e \in \mathbf{Int} \text{ or } (\exists
    e' \in \mathbf{Exp}.\ \exists \sigma' \in \mathbf{Store}.\ \langle \sigma, e \rangle \to
    \langle \sigma', e' \rangle)
    \]

    \item \textbf{Preservation:} Evaluation preserves containment of free variables in the
    domain of the store,
    \[
    \forall e, e' \in \mathbf{Exp}.\ \forall \sigma, \sigma' \in \mathbf{Store}.
    \]
    \[
    fvs(e) \subseteq \operatorname{dom}(\sigma) \text{ and } \langle \sigma, e \rangle \to
    \langle \sigma', e' \rangle \implies fvs(e') \subseteq \operatorname{dom}(\sigma').
    \]
\end{itemize}

The rest of this lecture shows how can we prove such properties using induction.

\section{Inductive sets}

Induction is an important concept in programming language theory. An inductively-defined set
$A$ is one that is described using a finite collection of axioms and inductive (inference)
rules. Axioms of the form
\[
a \in A
\]
indicate that $a$ is in the set $A$. Inductive rules
\[
\frac{a_1 \in A \quad \ldots \quad a_n \in A}{a \in A}
\]
indicate that if $a_1, \ldots, a_n$ are all elements of $A$, then $a$ is also an element of $A$.

The set $A$ is the set of all elements that can be inferred to belong to $A$ using a (finite)
number of applications of these rules, starting only from axioms. In other words, for each
element $a$ of $A$, we must be able to construct a finite proof tree whose final conclusion is
$a \in A$.

\begin{example}
The set described by a grammar is an inductive set. For instance, the set of arithmetic
expressions can be described with two axioms and three inference rules:
\[
\frac{}{x \in \mathbf{Exp}} \qquad \frac{}{n \in \mathbf{Exp}}
\]
\[
\frac{e_1 \in \mathbf{Exp} \quad e_2 \in \mathbf{Exp}}{e_1 + e_2 \in \mathbf{Exp}} \qquad
\frac{e_1 \in \mathbf{Exp} \quad e_2 \in \mathbf{Exp}}{e_1 * e_2 \in \mathbf{Exp}} \qquad
\frac{e_1 \in \mathbf{Exp} \quad e_2 \in \mathbf{Exp}}{x := e_1 ; e_2 \in \mathbf{Exp}}
\]
These axioms and rules describe the same set of expressions as the grammar:
\[
e ::= x \mid n \mid e_1 + e_2 \mid e_1 * e_2 \mid x := e_1 ; e_2
\]
\end{example}

\begin{example}
The natural numbers (expressed here in unary notation) can be inductively defined:
\[
\frac{}{0 \in \mathbb{N}} \qquad \frac{n \in \mathbb{N}}{\operatorname{succ}(n) \in \mathbb{N}}
\]
\end{example}

\begin{example}
The small-step evaluation relation $\to$ is an inductively defined set.
\end{example}

\begin{example}
The multi-step evaluation relation can be inductively defined:
\[
\frac{}{\langle \sigma, e \rangle \to^* \langle \sigma, e \rangle} \; \text{REFL}
\qquad
\frac{\langle \sigma, e \rangle \to \langle \sigma', e' \rangle \quad \langle \sigma', e'
\rangle \to^* \langle \sigma'', e'' \rangle}{\langle \sigma, e \rangle \to^* \langle \sigma'',
e'' \rangle} \; \text{TRANS}
\]
\end{example}

\begin{example}
The set of free variables of an expression $e$ can be inductively defined:
\[
\frac{}{y \in fvs(y)}
\]
\[
\frac{y \in fvs(e_1)}{y \in fvs(e_1 + e_2)} \qquad
\frac{y \in fvs(e_2)}{y \in fvs(e_1 + e_2)}
\]
\[
\frac{y \in fvs(e_1)}{y \in fvs(e_1 * e_2)} \qquad
\frac{y \in fvs(e_2)}{y \in fvs(e_1 * e_2)}
\]
\[
\frac{y \in fvs(e_1)}{y \in fvs(x := e_1 ; e_2)} \qquad
\frac{y \neq x \quad y \in fvs(e_2)}{y \in fvs(x := e_1 ; e_2)}
\]
\end{example}

\section{Inductive proofs}

We can prove facts about elements of an inductive set using an inductive reasoning that follows
the structure of the set definition.

\subsection{Mathematical induction}

You have probably seen proofs by induction over the natural numbers, called mathematical
induction. In such proofs, we typically want to prove that some property $P$ holds for all
natural numbers, that is, $\forall n \in \mathbb{N}.\ P(n)$. A proof by induction works by first
proving that $P(0)$ holds, and then proving for all $m \in \mathbb{N}$, if $P(m)$ then $P(m +
1)$. The principle of mathematical induction can be stated succinctly as
\[
P(0) \text{ and } (\forall m \in \mathbb{N}.\ P(m) \implies P(m+1)) \implies \forall n \in
\mathbb{N}.\ P(n).
\]
The proposition $P(0)$ is the basis of the induction (also called the base case) while $P(m)
\implies P(m+1)$ is called induction step (or the inductive case). While proving the induction
step, the assumption that $P(m)$ holds is called the induction hypothesis.

\subsection{Structural induction}

Given an inductively defined set $A$, to prove that a property $P$ holds for all elements of
$A$, we need to show:

\begin{enumerate}
    \item \textbf{Base cases:} For each axiom
    \[
    a \in A,
    \]
    $P(a)$ holds.

    \item \textbf{Inductive cases:} For each inference rule
    \[
    \frac{a_1 \in A \quad \ldots \quad a_n \in A}{a \in A},
    \]
    if $P(a_1)$ and $\ldots$ and $P(a_n)$ then $P(a)$.
\end{enumerate}

Note that if the set $A$ is the set of natural numbers from Example 2 above, then the
requirements for proving that $P$ holds for all elements of $A$ is equivalent to mathematical
induction.

If $A$ describes a syntactic set, then we refer to induction following the requirements above
as \emph{structural induction}. If $A$ is an operational semantics relation (such as the
small-step operational semantics relation $\to$) then such an induction is called
\emph{induction on derivations}. We will see examples of structural induction and induction on
derivations throughout the course.

\subsection{Example: Progress}

Let's consider the progress property defined above, and repeated here:

\medskip
\noindent\textbf{Progress:} For each store $\sigma$ and expression $e$ such that the free
variables of $e$ are contained in the domain of $\sigma$, either $e$ is an integer or there
exists a possible transition for $\langle \sigma, e \rangle$:
\[
\forall e \in \mathbf{Exp}.\ \forall \sigma \in \mathbf{Store}.\ fvs(e) \subseteq
\operatorname{dom}(\sigma) \implies e \in \mathbf{Int} \text{ or } (\exists e' \in
\mathbf{Exp}.\ \exists \sigma' \in \mathbf{Store}.\ \langle \sigma, e \rangle \to \langle
\sigma', e' \rangle)
\]

Let's rephrase this property in terms of an explicit predicate on expressions:
\[
P(e) \triangleq \forall \sigma \in \mathbf{Store}.\ fvs(e) \subseteq \operatorname{dom}(\sigma)
\implies e \in \mathbf{Int} \text{ or } (\exists e', \sigma'.\ \langle \sigma, e \rangle \to
\langle \sigma', e' \rangle)
\]

The idea is to build a proof that follows the inductive structure given by the grammar:
\[
e ::= x \mid n \mid e_1 + e_2 \mid e_1 * e_2 \mid x := e_1 ; e_2
\]
This technique is called ``structural induction on $e$.'' We analyze each case in the grammar
and show that $P(e)$ holds for that case. Since the grammar productions $e_1 + e_2$ and $e_1 *
e_2$ and $x := e_1 ; e_2$ are inductive, they are inductive steps in the proof; the cases for
$x$ and $n$ are base cases. The proof proceeds as follows.

\begin{proof}
Let $e$ be an expression. We will prove that
\[
\forall \sigma \in \mathbf{Store}.\ fvs(e) \subseteq \operatorname{dom}(\sigma) \implies e \in
\mathbf{Int} \text{ or } (\exists e', \sigma'.\ \langle \sigma, e \rangle \to \langle \sigma',
e' \rangle)
\]
by structural induction on $e$. We analyze several cases, one for each case in the grammar for
expressions:

\medskip
\noindent\textbf{Case $e = x$:} Let $\sigma$ be an arbitrary store, and assume that $fvs(e)
\subseteq \operatorname{dom}(\sigma)$. By the definition of $fvs$ we have $fvs(x) = \{x\}$. By
assumption we have $\{x\} \subseteq \operatorname{dom}(\sigma)$ and so $x \in
\operatorname{dom}(\sigma)$. Let $n = \sigma(x)$. By the \textsc{Var} axiom we have $\langle
\sigma, x \rangle \to \langle \sigma, n \rangle$, which finishes the case.

\medskip
\noindent\textbf{Case $e = n$:} We immediately have $e \in \mathbf{Int}$, which finishes the
case.

\medskip
\noindent\textbf{Case $e = e_1 + e_2$:} Let $\sigma$ be an arbitrary store, and assume that
$fvs(e) \subseteq \operatorname{dom}(\sigma)$. We will assume that $P(e_1)$ and $P(e_2)$ hold
and show that $P(e)$ holds. Let's expand these properties. We have
\[
P(e_1) = \forall \sigma \in \mathbf{Store}.\ fvs(e_1) \subseteq \operatorname{dom}(\sigma)
\implies e_1 \in \mathbf{Int} \text{ or } (\exists e', \sigma'.\ \langle \sigma, e_1 \rangle \to
\langle \sigma', e' \rangle)
\]
\[
P(e_2) = \forall \sigma \in \mathbf{Store}.\ fvs(e_2) \subseteq \operatorname{dom}(\sigma)
\implies e_2 \in \mathbf{Int} \text{ or } (\exists e', \sigma'.\ \langle \sigma, e_2 \rangle \to
\langle \sigma', e' \rangle)
\]
and want to prove:
\[
P(e_1 + e_2) = \forall \sigma \in \mathbf{Store}.\ fvs(e_1+e_2) \subseteq
\operatorname{dom}(\sigma) \implies e_1 + e_2 \in \mathbf{Int} \text{ or } (\exists e',
\sigma'.\ \langle \sigma, e_1 + e_2 \rangle \to \langle \sigma', e' \rangle)
\]
We analyze several subcases.

\smallskip
\noindent\textit{Subcase $e_1 = n_1$ and $e_2 = n_2$:} By rule \textsc{Add}, we immediately have
$\langle \sigma, n_1 + n_2 \rangle \to \langle \sigma, p \rangle$, where $p = n_1 + n_2$.

\smallskip
\noindent\textit{Subcase $e_1 \notin \mathbf{Int}$:} By assumption and the definition of $fvs$
we have
\[
fvs(e_1) \subseteq fvs(e_1 + e_2) \subseteq \operatorname{dom}(\sigma)
\]
Hence, by the induction hypothesis $P(e_1)$ we also have $\langle \sigma, e_1 \rangle \to
\langle \sigma', e' \rangle$ for some $e'$ and $\sigma'$. By rule \textsc{Ladd} we have $\langle
\sigma, e_1 + e_2 \rangle \to \langle \sigma', e' + e_2 \rangle$.

\smallskip
\noindent\textit{Subcase $e_1 = n_1$ and $e_2 \notin \mathbf{Int}$:} By assumption and the
definition of $fvs$ we have
\[
fvs(e_2) \subseteq fvs(e_1 + e_2) \subseteq \operatorname{dom}(\sigma)
\]
Hence, by the induction hypothesis $P(e_2)$ we also have $\langle \sigma, e_2 \rangle \to
\langle \sigma', e' \rangle$ for some $e'$ and $\sigma'$. By rule \textsc{Radd} we have $\langle
\sigma, e_1 + e_2 \rangle \to \langle \sigma', e_1 + e' \rangle$, which finishes the case.

\medskip
\noindent\textbf{Case $e = e_1 * e_2$:} Analogous to the previous case.

\medskip
\noindent\textbf{Case $e = x := e_1 ; e_2$:} Let $\sigma$ be an arbitrary store, and assume that
$fvs(e) \subseteq \operatorname{dom}(\sigma)$. As above, we assume that $P(e_1)$ and $P(e_2)$
hold and show that $P(e)$ holds. Let's expand these properties. We have
\[
P(e_1) = \forall \sigma.\ fvs(e_1) \subseteq \operatorname{dom}(\sigma) \implies e_1 \in
\mathbf{Int} \text{ or } (\exists e', \sigma'.\ \langle \sigma, e_1 \rangle \to \langle \sigma',
e' \rangle)
\]
\[
P(e_2) = \forall \sigma.\ fvs(e_2) \subseteq \operatorname{dom}(\sigma) \implies e_2 \in
\mathbf{Int} \text{ or } (\exists e', \sigma'.\ \langle \sigma, e_2 \rangle \to \langle \sigma',
e' \rangle)
\]
and want to prove:
\[
P(x := e_1 ; e_2) = \forall \sigma.\ fvs(x := e_1 ; e_2) \subseteq \operatorname{dom}(\sigma)
\implies
\]
\[
x := e_1 ; e_2 \in \mathbf{Int} \text{ or } (\exists e', \sigma'.\ \langle \sigma, x := e_1 ;
e_2 \rangle \to \langle \sigma', e' \rangle)
\]
We analyze several subcases.

\smallskip
\noindent\textit{Subcase $e_1 = n_1$:} By rule \textsc{Assgn} we have $\langle \sigma, x := n_1
; e_2 \rangle \to \langle \sigma', e_2 \rangle$ where $\sigma' = \sigma[x \mapsto n_1]$.

\smallskip
\noindent\textit{Subcase $e_1 \notin \mathbf{Int}$:} By assumption and the definition of $fvs$
we have
\[
fvs(e_1) \subseteq fvs(x := e_1 ; e_2) \subseteq \operatorname{dom}(\sigma)
\]
Hence, by induction hypothesis we also have $\langle \sigma, e_1 \rangle \to \langle \sigma',
e' \rangle$ for some $e'$ and $\sigma'$. By the rule \textsc{Assgn1} we have $\langle \sigma, x
:= e_1 ; e_2 \rangle \to \langle \sigma', x := e_1' ; e_2 \rangle$, which finishes the case and
the inductive proof.
\end{proof}

\end{document}
